\chapter*{Exposé XXVI\\Parabolic Subgroups of Reductive Groups}
\addcontentsline{toc}{chapter}{Exposé XXVI. Parabolic Subgroups of Reductive Groups}
\begin{center}
\large M. Demazure
\end{center}

\medskip

This Exposé studies the parabolic subgroups of an \(S\)-reductive group \(G\).
Its essential result is the conjugacy theorem (\SGARefLink{sga3:xxvi:cor:5:4}{5.4}).  Its essential tool is
the notion of the transversal position of two parabolic subgroups, a notion
studied systematically in §4.  Another fact plays an important role: the
decomposition of the unipotent radical \(\rad^{u}(P)\) of a parabolic subgroup
\(P\) into successive extensions of vector groups
(\SGARefLink{sga3:xxvi:proposition:2-1}{2.1}).\footnote{\emph{Editor's note.} Since the Levi subgroups of \(P\) form a
torsor under \(\rad^{u}(P)\) (\SGARefLink{sga3:xxvi:cor:1:9}{1.9}), this entails, when \(S\) is semi-local,
that \(P\) has a Levi subgroup and hence a maximal torus (\SGARefLink{sga3:xxvi:cor:2:4}{2.4}).}

Various schemes associated with \(G\) are studied in §3; §6 deals with the
split subtori\footnote{\emph{Editor's note.} The terminology ``trivial
torus'' has been replaced by ``split torus''.} of \(G\) and their relations
with parabolic subgroups.

Finally, in §7, we explain briefly how, over a semi-local base, one formulates
the ``relative theory'' of reductive groups as presented over fields in the
article of A.~Borel and J.~Tits, \emph{Groupes réductifs}, Publications
Mathématiques de l'IHÉS, no.~27.  In that article, cited as [BT65] below, the
reader will also find, for a base field, other results not touched upon here.
\footnote{\emph{Editor's note.} Version of 13 October 2024.}

\SGASection{1}{Recollections. Levi subgroups}
\SGARefTarget{sga3:xxvi:section:1}

\begin{sgastatement}{Definition 1.1}
Let \(S\) be a scheme, \(G\) an \(S\)-reductive group, and \(P\) a
sub-\(S\)-group-scheme of \(G\).  One says that \(P\) is a
\emph{parabolic subgroup} of \(G\) if:
\begin{enumerate}[label=\textup{(\roman*)}]
\item \(P\) is smooth over \(S\);
\item for every \(s\in S\), the quotient \(\bar s\)-scheme
\(G_{\bar s}/P_{\bar s}\) is proper (that is, by the Bible, §\SGARefLink{sga3:xxvi:corollary:6-4}{6.4},
\SGARefLink{sga3:xxvi:prop:1:4}{Theorem~4} \(=\) [Ch05], §\SGARefLink{sga3:xxvi:subsec:6:5}{6.5}, Theorem~5, \(P_{\bar s}\) contains a
Borel subgroup of \(G_{\bar s}\)).
\end{enumerate}
\end{sgastatement}
\SGARefTarget{sga3:xxvi:def:1:1}

\begin{sgastatement}{Proposition 1.2 \textup{(Exp.\ XXII, \SGARefLink{sga3:xxii:cor:5.8.5}{5.8.5})}}
Let \(S\) be a scheme, \(G\) an \(S\)-reductive group, and \(P\) a parabolic
subgroup of \(G\).  Then \(P\) is closed in \(G\), has connected fibers, and
\[
  P=\Norm_G(P).
\]
Moreover, the quotient sheaf \(G/P\) is representable by an \(S\)-scheme
that is smooth and projective over \(S\).
\end{sgastatement}
\SGARefTarget{sga3:xxvi:prop:1:2}

\begin{sgastatement}{\SGAInlineRefTarget{sga3:xxvi:proposition:1-3}Proposition 1.3 \textup{(Exp.\ XXII, \SGARefLink{sga3:xxii:thm:5.3.9}{5.3.9} and 5.3.11)}}
Let \(S\) be a scheme, \(G\) an \(S\)-reductive group, and \(P,P'\) two
parabolic subgroups of \(G\).  The following conditions are equivalent:
\begin{enumerate}[label=\textup{(\roman*)}]
\item \(P\) and \(P'\) are conjugate in \(G\), locally for the étale
\textup{(respectively, fpqc)} topology.
\item For every \(s\in S\), \(P_{\bar s}\) and \(P'_{\bar s}\) are conjugate
by an element of \(G(\bar s)\).
\item The strict transporter \(\Transt_G(P,P')\) of \(P\) into \(P'\), defined
for every \(S'\to S\) by
\[
 \Transt_G(P,P')(S')
 =\{g\in G(S')\mid \Int(g)P_{S'}=P'_{S'}\},
\]
is a closed subscheme of \(G\), smooth and of finite presentation over
\(S\), and is a principal homogeneous bundle on the right under \(P\) and
on the left under \(P'\).
\end{enumerate}
\end{sgastatement}
\SGARefTarget{sga3:xxvi:prop:1:3}

\begin{sgastatement}{Proposition 1.4}
Let \(S\) be a nonempty scheme, let \((G,T,M,R)\) be a split
\(S\)-reductive group, and let \(R'\) be a subset of \(R\).  The following
conditions on \(R'\) are equivalent:
\begin{enumerate}[label=\textup{(\roman*)}]
\item
\[
 \mathfrak g_{R'}=\mathfrak t\oplus
 \bigoplus_{\alpha\in R'}\mathfrak g^\alpha
\]
is the Lie algebra of a parabolic subgroup of \(G\) containing \(T\)
\textup{(necessarily unique, Exp.\ XXII, \SGARefLink{sga3:xxii:cor:5.3.5}{5.3.5})}.
\item \(R'\) is of type \textup{(R)} \textup{(Exp.\ XXII, \SGARefLink{sga3:xxii:def:5.4.2}{5.4.2})} and
contains a system of positive roots.
\item \(R'\) is a closed subset of \(R\) and satisfies: if
\(\alpha\in R-R'\), then \(-\alpha\in R'\); equivalently,
\(R=R'\cup(-R')\).
\item There exist a system \(\Delta\) of simple roots and a subset
\(A\subset\Delta\) such that \(R'\) is the union of the positive roots and
the negative roots that are linear combinations of elements of \(A\).
\item\SGAInlineRefTarget{sga3:xxvi:prop:1:4:item:v} \(R'\) contains a system of simple roots of \(R\); moreover, if
\(\Delta\subset R'\) is a system of simple roots of \(R\) and
\[
 A=(-R')\cap\Delta,
\]
then \(R'\) is the union of the positive roots and the negative roots that
are linear combinations of elements of \(A\).
\end{enumerate}
\end{sgastatement}
\SGARefTarget{sga3:xxvi:prop:1:4}

One has (i)\(\Leftrightarrow\)(ii) by Exp.\ XXII, \SGARefLink{sga3:xxii:prop:5.4.5:item:ii}{5.4.5(ii)} and \SGARefLink{sga3:xxii:prop:5.5.1}{5.5.1},
and (iii)\(\Rightarrow\)(ii) by Exp.\ XXI, \SGARefLink{sga3:XXI:3.3.6}{3.3.6} and Exp.\ XXII, \SGARefLink{sga3:xxii:thm:5.4.7}{5.4.7}.
Clearly (v)\(\Rightarrow\)(iv)\(\Rightarrow\)(iii), while
(iii)\(\Rightarrow\)(v) follows from Exp.\ XXI, \SGARefLink{sga3:XXI:3.3.6}{3.3.6} and 3.3.10.

It remains to prove that (i) implies that \(R'\) is closed in \(R\).  This
last assertion may be checked on any geometric fiber, so we may suppose that
\(S\) is the spectrum of an algebraically closed field.

Let \(P\) be the parabolic subgroup of \(G\) containing \(T\) whose Lie
algebra is \(\mathfrak g_{R'}\).  Since the Borel subgroups of \(P\) are the
Borel subgroups of \(G\) contained in \(P\), the Bible, §12.3, \SGARefLink{sga3:xxvi:theorem:4:2:1}{Theorem~1},
and Exp.\ XXII, \SGARefLink{sga3:xxii:prop:5.4.5:item:i}{5.4.5(i)}, show that, if \(U\) denotes the unipotent radical
of \(P\), then \(T\cdot U\) is the subgroup of \(G\) containing \(T\) whose
Lie algebra is
\[
 \mathfrak g_{R''}=\mathfrak t\oplus
 \bigoplus_{\alpha\in R''}\mathfrak g^\alpha,
\]
where \(R''\) is the intersection of the systems of positive roots of \(R\)
contained in \(R'\).  In particular, \(R''\) is closed and
\(R''\cap(-R'')=\varnothing\).

On the other hand, \(H=P/U\) is reductive, the canonical image \(\bar T\)
of \(T\) is a maximal torus of \(H\) (and \(T\to\bar T\) is an
isomorphism), and there is an isomorphism of \(\bar T\)-modules, or
equivalently of \(M\)-graded vector spaces,
\[
 \operatorname{Lie}(H)\simeq\mathfrak t\oplus
 \bigoplus_{\alpha\in R_s}\mathfrak g^\alpha,
\]
where \(R_s\) is the complement of \(R''\) in \(R'\).  Thus \(R_s\)
identifies naturally with the roots of \(H\) relative to \(\bar T\), and in
particular \(R_s=-R_s\).  Consequently,
\[
 R''=\{\alpha\in R'\mid-\alpha\notin R'\},
 \qquad
 R_s=\{\alpha\in R'\mid-\alpha\in R'\}.
\]

We now prove that \(R'\) is closed.  Let \(\alpha,\beta\in R'\) with
\(\alpha+\beta\in R\).  If \(\alpha,\beta\in R''\), then
\(\alpha+\beta\in R''\) because \(R''\) is closed.  If
\(\alpha\in R_s\), \(\beta\in R''\), and \(\alpha+\beta\notin R'\), then
\(\alpha+\beta\in-R''\), whence
\[
 -\alpha=-(\alpha+\beta)+\beta\in R'',
\]
because \(R''\) is closed.  This gives
\(-\alpha\in R''\cap R_s\), a contradiction.  It remains to consider
\(\alpha,\beta\in R_s\).  If \(\alpha+\beta\notin R'\), then
\(\alpha+\beta\in-R''\).  Since \(\alpha+\beta\ne0\), some system of
positive roots of the root system \(R_s\) contains both \(\alpha\) and
\(\beta\), and hence some Borel subgroup of \(H=P/U\) contains the canonical
images of \(U_\alpha\) and \(U_\beta\).  Its inverse image in \(P\) is a
Borel subgroup containing \(U_\alpha\), \(U_\beta\), \(U\), and therefore
\(U_{-(\alpha+\beta)}\), which is impossible.

\begin{sgastatement}{Corollary 1.5}
A parabolic subgroup of a reductive group is of type \textup{(RC)}
\textup{(Exp.\ XXII, \SGARefLink{sga3:xxii:def:5.11.1}{5.11.1})}.
\end{sgastatement}
\SGARefTarget{sga3:xxvi:cor:1:5}

\begin{sgastatement}{\SGAInlineRefTarget{sga3:xxvi:proposition:1-6}Proposition 1.6 \textup{(Exp.\ XXII, \SGARefLink{sga3:xxii:prop:5.11.4}{5.11.4})}}
Let \(S\) be a scheme, \(G\) an \(S\)-reductive group, and \(P\) a subgroup
of \(G\) of type \textup{(RC)}.\footnote{\emph{Editor's note.}
``Parabolic subgroup'' has been replaced by ``subgroup of type (RC)'', as
in Exp.\ XXII, \SGARefLink{sga3:xxii:prop:5.11.4}{5.11.4}, because the statement will later be applied
(\SGARefLink{sga3:xxvi:proposition:4-5-1}{4.5.1}, \SGARefLink{sga3:xxvi:cor:6:17}{6.17}) to \(P\cap P'\) when \(P,P'\) are parabolic and
\(P\cap P'\) is of type (RC).}
\begin{enumerate}[label=\textup{(\roman*)}]
\item \(P\) has a largest invariant sub-group-scheme that is smooth and of
finite presentation over \(S\), with connected unipotent geometric fibers.
It is a characteristic subgroup of \(P\), called the \emph{unipotent
radical} of \(P\) and denoted \(\rad^u(P)\).  The quotient sheaf
\(P/\rad^u(P)\) is representable by an \(S\)-reductive group.
\item\SGAInlineRefTarget{sga3:xxvi:proposition:1-6:item:ii} If \(T\) is a maximal torus of \(P\), then \(P\) has a reductive
subgroup \(L\) containing \(T\) such that:
\begin{enumerate}[label=\textup{(\alph*)}]
\item every reductive subgroup of \(P\) containing \(T\) is contained in
\(L\);
\item \(P\) is the semidirect product \(L\cdot\rad^u(P)\), that is, the
canonical morphism \(L\to P/\rad^u(P)\) is an isomorphism.
\end{enumerate}
Moreover, \(L\) is the unique subgroup \textup{(respectively, reductive
subgroup)} of \(P\) containing \(T\) and satisfying (b) \textup{(respectively,
(a))}.  Finally,
\[
 \Norm_P(L)=L,\qquad \Norm_P(T)=\Norm_L(T).
\]
\end{enumerate}
\end{sgastatement}
\SGARefTarget{sga3:xxvi:prop:1:6}

\paragraph{\SGARefLink{sga3:xxvi:prop:1:6}{1.7}.}
A subgroup \(L\) of \(P\) satisfying condition \SGARefLink{sga3:xxvi:proposition:1-6:item:ii}{1.6(ii)}(b) is called a
\emph{Levi subgroup} of \(P\).  It is a maximal reductive subgroup of
\(P\).  Indeed, it is reductive because it is isomorphic to
\(P/\rad^u(P)\).  If \(L'\) is a reductive subgroup of \(P\) containing
\(L\), then, locally for the fpqc topology, \(L\) has a maximal torus
\(T\), and \SGARefLink{sga3:xxvi:proposition:1-6:item:ii}{1.6(ii)} gives \(L'=L\).
\SGARefTarget{sga3:xxvi:item:1:7}

If \(L\) and \(L'\) are two Levi subgroups of \(P\), they are conjugate in
\(P\), locally for the fpqc topology.  Indeed, locally one may choose
maximal tori \(T\subset L\) and \(T'\subset L'\); after a further fpqc
localization they are conjugate in \(P\), so one may assume \(T=T'\), and
then \SGARefLink{sga3:xxvi:proposition:1-6:item:ii}{1.6(ii)} gives \(L=L'\).  Since
\(P=L\cdot\rad^u(P)\) and \(\Norm_P(L)=L\), one obtains:

\begin{sgastatement}{Corollary 1.8}
Let \(P\) be a subgroup of type \textup{(RC)} of the \(S\)-reductive group
\(G\).  If \(L,L'\) are two Levi subgroups of \(P\), there is a unique
\(u\in\rad^u(P)(S)\) such that \(\Int(u)L=L'\).
\end{sgastatement}
\SGARefTarget{sga3:xxvi:cor:1:8}

Let \(\Lev(P)\) denote the functor of Levi subgroups of \(P\): for
\(S'\to S\), \(\Lev(P)(S')\) is the set of Levi subgroups of \(P_{S'}\).
\SGARefLink{sga3:xxvi:cor:1:8}{Corollary~1.8} gives:

\begin{sgastatement}{Corollary 1.9}
Let \(P\) be a subgroup of type \textup{(RC)} of the \(S\)-reductive group
\(G\).  Then \(\Lev(P)\) is a principal homogeneous bundle under the
\(S\)-group \(\rad^u(P)\); in particular, it is representable by an
\(S\)-scheme that is smooth and affine over \(S\), with integral geometric
fibers.
\end{sgastatement}
\SGARefTarget{sga3:xxvi:cor:1:9}

Likewise, \SGARefLink{sga3:xxvi:proposition:1-6}{1.6} immediately yields:

\begin{sgastatement}{Corollary 1.10}
Let \(P\) be a subgroup of type \textup{(RC)} of the \(S\)-reductive group
\(G\).  The functor \(\Tor(P)\) of maximal tori of \(P\) is representable
by a smooth affine \(S\)-scheme.  The relation \(L\supset T\) defines a
morphism
\[
 \Tor(P)\longrightarrow\Lev(P),
\]
whose fiber above \(L\in\Lev(P)(S)\) identifies with \(\Tor(L)\)
\textup{(Exp.\ XXII, \SGARefLink{sga3:xxii:cor:5.8.3}{5.8.3})}.
\end{sgastatement}
\SGARefTarget{sga3:xxvi:cor:1:10}

The first assertion of \SGARefLink{sga3:xxvi:cor:1:10}{1.10} follows from the other two and Exp.\ XXII,
\SGARefLink{sga3:xxii:cor:5.8.3}{5.8.3}.

\begin{sgastatement}{\SGAInlineRefTarget{sga3:xxvi:definition:1-11}Definition 1.11}
Let \(S\) be a nonempty scheme, \(G\) an \(S\)-reductive group, \(P\) a
parabolic subgroup of \(G\), and
\[
 \mathcal E=(T,M,R,\Delta,(X_\alpha)_{\alpha\in\Delta})
\]
a pinning of \(G\).  One says that \(\mathcal E\) is \emph{adapted} to
\(P\), or is a \emph{pinning of the pair \((G,P)\)}, if \(P\supset T\) and
the Lie algebra of \(P\) is of the form
\[
 \mathfrak t\oplus\bigoplus_{\alpha\in R'}\mathfrak g^\alpha,
\]
where \(R'\subset R\) contains \(R^+\).
\end{sgastatement}
\SGARefTarget{sga3:xxvi:def:1:11}

In particular, if \(T\subset B\) is the Killing pair defined by the
pinning, then \(T\subset B\subset P\).  Under these conditions set
\(\Delta(P)=\Delta\cap(-R')\).  Exp.\ XXII, \SGARefLink{sga3:xxii:lem:5.4.3}{5.4.3} gives
\[
 \alpha\in\Delta(P)
 \ \Longleftrightarrow\
 \alpha\in\Delta\text{ and }U_{-\alpha}\subset P
 \ \Longleftrightarrow\
 \alpha\in\Delta\text{ and }U_{-\alpha}\cap P\ne e.
\]

\SGARefLink{sga3:xxvi:prop:1:4:item:v}{Proposition~1.4(v)} and Exp.\ XXII, \SGARefLink{sga3:xxii:lem:5.11.3}{5.11.3} and 5.10.6 now immediately give:

\begin{sgastatement}{\SGAInlineRefTarget{sga3:xxvi:proposition:1-12}Proposition 1.12}
Let \(S\) be a scheme, \(G\) an \(S\)-reductive group, \(P\) a parabolic
subgroup of \(G\), and
\((T,M,R,\Delta,(X_\alpha)_{\alpha\in\Delta})\) a pinning of \(G\) adapted
to \(P\).  Let \(\Delta(P)\subset\Delta\) be as above.
\begin{enumerate}[label=\textup{(\roman*)}]
\item The unipotent radical of \(P\) is
\[
 \rad^u(P)=U_{R''}=\prod_{\alpha\in R''}U_\alpha,
\]
where \(R''\) is the set of positive roots whose expansion in \(\Delta\)
contains at least one element of
\(\Delta-\Delta(P)\)\footnote{\emph{Editor's note.}
\(\Delta(P)\) has been corrected to \(\Delta-\Delta(P)\).} with a
nonzero coefficient.
\item The unique Levi subgroup \(L\) of \(P\) containing \(T\) is
\[
 Z_{\Delta(P)}=\Centr_G(T_{\Delta(P)}),
\]
where \(T_{\Delta(P)}\) is the maximal torus of
\(\bigcap_{\alpha\in\Delta(P)}\ker(\alpha)\).  Moreover,
\(T_{\Delta(P)}=\rad(L)\).
\end{enumerate}
\end{sgastatement}
\SGARefTarget{sga3:xxvi:prop:1:12}

\begin{sgastatement}{Corollary 1.13}
Every Levi subgroup \(L\) of a parabolic subgroup \(P\) of a reductive group
\(G\) is a \emph{critical subgroup} of \(G\); that is
\textup{(Exp.\ XXII, \SGARefLink{sga3:xxii:def:5.10.4}{5.10.4})},
\[
 L=\Centr_G(\rad(L)).
\]
\end{sgastatement}
\SGARefTarget{sga3:xxvi:cor:1:13}

This follows immediately from \SGARefLink{sga3:xxvi:proposition:1-12}{1.12} and the following lemma, which is
contained in \SGARefLink{sga3:xxvi:prop:1:3}{1.4} and Exp.\ XXII, \SGARefLink{sga3:xxii:5.4.1}{5.4.1}:

\begin{sgastatement}{Lemma 1.14}
Locally for the étale topology, every pair \((G,P)\), where \(P\) is a
parabolic subgroup of the reductive group \(G\), admits a pinning in the
sense of \SGARefLink{sga3:xxvi:definition:1-11}{1.11}.
\end{sgastatement}
\SGARefTarget{sga3:xxvi:lemma:1:14}

We also record:

\begin{sgastatement}{Proposition 1.15}
Let \(S\) be a scheme, \(G\) an \(S\)-reductive group, \(P\) a parabolic
subgroup of \(G\), and \(\mathcal E,\mathcal E'\) two pinnings of \(G\)
adapted to \(P\).  The unique inner automorphism of \(G\) over \(S\) that
transforms \(\mathcal E\) into \(\mathcal E'\) \textup{(Exp.\ XXIV, \SGARefLink{sga3:xxiv:1.5}{1.5})}
comes from \(P\), via
\[
 P\longrightarrow P/\Centr(P)
 =P/\Centr(G)\longrightarrow G/\Centr(G).
\]
\end{sgastatement}
\SGARefTarget{sga3:xxvi:prop:1:15}

Indeed, it suffices to argue as in Exp.\ XXIV, \SGARefLink{sga3:xxiv:1.5}{1.5}, using:

\begin{sgastatement}{Lemma 1.16 \textup{(Exp.\ XXII, \SGARefLink{sga3:xxii:cor:5.3.14}{5.3.14} and \SGARefLink{sga3:xxii:prop:5.2.6}{5.2.6})}}
The maximal tori \textup{(respectively, Borel subgroups, respectively,
Killing pairs)} of a parabolic subgroup \(P\) of an \(S\)-reductive group
\(G\) are conjugate in \(P\), locally for the étale topology.
\end{sgastatement}
\SGARefTarget{sga3:xxvi:lemma:1:16}

\begin{sgastatement}{Proposition 1.17}
Let \(S\) be a scheme, \(G\) an \(S\)-reductive group, \(P,P'\) two
parabolic subgroups of \(G\), and \(B\) a Borel subgroup contained in both
\(P\) and \(P'\).  If \(P\) and \(P'\) are conjugate in \(G\), locally for
the étale topology, then \(P=P'\).
\end{sgastatement}
\SGARefTarget{sga3:xxvi:prop:1:17}

Indeed, we may assume that \(g\in G(S)\) satisfies \(\Int(g)P=P'\).  The
groups \(B\) and \(\Int(g)^{-1}B\) are then two Borel subgroups of \(P\).
After extending \(S\), \SGARefLink{sga3:xxvi:lemma:1:16}{Lemma~1.16} lets us choose \(p\in P(S)\) with
\(\Int(p)\Int(g^{-1})B=B\).  Hence
\[
 pg^{-1}\in\Norm_G(B)(S)=B(S),
\]
so \(g\in B(S)\cdot p\subset P(S)\), and therefore
\(P'=\Int(g)P=P\).

\begin{sgastatement}{Remark 1.18}
If \(P,P'\) are two parabolic subgroups of \(G\) containing a common Borel
subgroup, then \(P\cap P'\) is again a parabolic subgroup of \(G\).  Indeed,
it is smooth along the unit section \textup{(Exp.\ XXII, \SGARefLink{sga3:xxii:prop:5.4.5}{5.4.5})} and
contains a Borel subgroup.
\end{sgastatement}
\SGARefTarget{sga3:xxvi:rem:1:18}

\begin{sgastatement}{\SGAInlineRefTarget{sga3:xxvi:proposition:1-19}Proposition 1.19}
Let \(S\) be a scheme, \(G\) an \(S\)-reductive group, and \(G'\) its
derived group \textup{(Exp.\ XXII, \SGARefLink{sga3:xxii:thm:6.2.1}{6.2.1})}.
\begin{enumerate}[label=\textup{(\roman*)}]
\item The assignments
\[
 P\longmapsto P'=P\cap G',
 \qquad
 P'\longmapsto P'\cdot\rad(G)=\Norm_G(P')
\]
are mutually inverse bijections between the parabolic subgroups of \(G\)
and those of \(G'\).  Moreover,
\(\rad^u(P)=\rad^u(P')\).
\item Let \(P\) be a parabolic subgroup of \(G\) and \(P'=P\cap G'\).
The assignments
\[
 L\longmapsto L'=L\cap G'=L\cap P',
 \qquad
 L'\longmapsto L'\cdot\rad(G)=\Centr_G(\rad(L'))
\]
are mutually inverse bijections between the Levi subgroups of \(P\) and
those of \(P'\).  Moreover,
\[
 \rad(L')=(\rad(L)\cap G')^0.
\]
\end{enumerate}
\end{sgastatement}
\SGARefTarget{sga3:xxvi:prop:1:19}

The proof, for example by reduction to the split case, is straightforward
and is left to the reader, together with the immediate proof of:

\begin{sgastatement}{\SGAInlineRefTarget{sga3:xxvi:proposition:1-20}Proposition 1.20}
Let \(S\) be a scheme, \(G\) an \(S\)-reductive group, \(P\) a parabolic
subgroup of \(G\), and \(L\) a Levi subgroup of \(P\).  The assignments
\[
 Q\longmapsto Q\cap L=Q',
 \qquad
 Q'\longmapsto Q'\cdot\rad^u(P)
\]
are mutually inverse bijections between the parabolic subgroups of \(G\)
contained in \(P\) and the parabolic subgroups of \(L\).  Moreover, the
Levi subgroups of \(Q'\) are precisely the Levi subgroups of \(Q\)
contained in \(L\).
\end{sgastatement}
\SGARefTarget{sga3:xxvi:prop:1:20}

\SGARefLink{sga3:xxvi:proposition:1-6}{Proposition~1.6} may be completed as follows.

\begin{sgastatement}{\SGAInlineRefTarget{sga3:xxvi:proposition:1-21}Proposition 1.21}
Let \(S\) be a scheme, \(G\) an \(S\)-reductive group, and \(P\) a
parabolic subgroup of \(G\).
\begin{enumerate}[label=\textup{(\roman*)}]
\item \(P\) has a largest invariant subgroup that is smooth and of finite
presentation over \(S\), with connected solvable geometric fibers.  It is
a characteristic subgroup of \(P\), called the \emph{radical} of \(P\)
and denoted \(\rad(P)\).  The quotient sheaf \(P/\rad(P)\) is representable
by an \(S\)-semisimple group.
\item If \(L\) is a Levi subgroup of \(P\), then \(\rad(P)\) is the
semidirect product of \(\rad^u(P)\) and \(\rad(L)\); one has
\[
 \rad(L)=L\cap\rad(P),
 \qquad
 L=\Centr_G(L\cap\rad(P)),
 \qquad
 P/\rad(P)\simeq L/\rad(L).
\]
\end{enumerate}
\end{sgastatement}
\SGARefTarget{sga3:xxvi:prop:1:21}

Indeed, assertion (i) being local, one may suppose that \(P\) has a Levi
subgroup \(L\).  It is then immediate that
\[
 R=\rad^u(P)\cdot\rad(L)
\]
has the properties asserted in (i).  For (ii), it only remains to prove
that \(\rad(L)=L\cap\rad(P)\); this follows at once because
\(L\cap\rad(P)\) is smooth with connected fibers, \(L\) being the
centralizer of a torus.\footnote{\emph{Editor's note.} See \SGARefLink{sga3:xix:lemma:1:4}{Exposé XIX,
1.4}.}
